% docs/texte/20_core/observation.tex -- chapter "Observation" (Interferosim)
% Version: 0.2.0 (2026-10-31) -- section "Geometry" (module observation.geometry)
%   0.1.0 (2026-10-12): placeholder
% ================================================================
\label{ch:observation}

Hour angle and elevation (sidereal and solar day), uv-coverage with
Earth rotation synthesis, visibilities with noise, adding and
multiplying interferometer, fringe rate and delay.
The uv-coverage, the visibility simulation with noise and the transit
follow with version~0.5; the geometry needed by the signal simulation
(chapter~\ref{ch:signal}) is implemented since version~0.4.

\section{Geometry}
\index{geometric delay}\index{uvw}

The module \texttt{interferosim.observation.geometry} computes the source
direction in the earth-fixed frame (ECEF, chapter~\ref{ch:array}) from
right ascension $\alpha$, declination $\delta$ and the Greenwich hour angle
$H_G = \mathrm{GMST} - \alpha$ (or $H_G = H - \lambda$ from the local hour
angle $H$ at a site of longitude $\lambda$):
\begin{align}
  \hat s_0 &= (\cos\delta\cos H_G,\; -\cos\delta\sin H_G,\; \sin\delta),\\
  \hat e_E &= (\sin H_G,\; \cos H_G,\; 0),\qquad
  \hat e_N = (-\sin\delta\cos H_G,\; \sin\delta\sin H_G,\; \cos\delta),
\end{align}
$\hat s_0$ towards the phase centre, $\hat e_E$ and $\hat e_N$ east
(increasing $\alpha$) and north on the sky. A source at the offset $(l, m)$
lies in the direction $\hat s = \sqrt{1-l^2-m^2}\,\hat s_0 + l\,\hat e_E +
m\,\hat e_N$. GMST follows the IAU~1982 formula
$\mathrm{GMST} = 280.46061837^\circ + 360.98564736629^\circ\,(JD-2451545)$;
precession and nutation are neglected (AAS6 uses astropy; for the simulation
it suffices that simulator and correlator use the same model).

For a baseline $\vec B = \vec r_j - \vec r_i$ (station $i$ the reference):
\begin{align}
  \tau_g &= -\frac{\vec B\cdot\hat s}{c}
     && \texttt{geometric\_delay}\\
  (u, v, w) &= \frac{1}{\lambda}\,\bigl(\vec B\cdot\hat e_E,\;
               \vec B\cdot\hat e_N,\; \vec B\cdot\hat s_0\bigr)
     && \texttt{uvw}
\end{align}
$\tau_g > 0$ means the wavefront reaches station $j$ later. The delay rate
\texttt{delay\_rate} follows from the Earth's rotation
($\dot H_G = \omega_E$); for an east--west baseline it is
$\omega_E\,b\cos\delta\cos H/c$, maximal at transit. \texttt{elevation\_deg}
gives the elevation of a direction at a site.
